\begin{frame}
  \frametitle{Question 4: 辞書式順序の利点}
  \textbf{問題:} この教科書で定義された辞書式順序 (the lexicographic order) は，いわゆる辞書順（辞書に単語が掲載される順）とは異なる．例えば，辞書では appropriateは backより前に掲載されるが，辞書式順序では back<appropriateである（短い文字列が長い文字列より前）．文字列の順序をこのように定義する理由（目的や利点）を考えて（または想像して）答えなさい．
  \vspace{0.3em}

  \begin{alertblock}{解答1}
      辞書式順序は、アルファベット順だけでなく、文字列の構造も考慮するため、
      特定のアルゴリズムやデータ構造、例えば深さ優先探索において、より自然な順序を提供することができる。
      辞書式順序だと文字列が短いものから長いものになるのでコンピュータのソートが簡単になるので処理速度が上がる
  \end{alertblock}
\end{frame}